\documentclass[10pt,a4paper]{amsart}
\usepackage[margin=19mm]{geometry}
\usepackage{amsmath,amssymb,mathtools}
\usepackage[scr=rsfs,cal=euler]{mathalfa}
\usepackage{xfrac}
\usepackage[unicode,hidelinks,bookmarks=false]{hyperref}
\newcommand{\ie}{i.e.\ }
\newcommand{\cf}{cf.\ }
\newcommand{\resp}{respectively\ }
\newcommand{\Subsection}{\subsection}
\providecommand{\nobreakdash}{\nobreak\hbox{-}}
\setlength{\emergencystretch}{2em}
\multlinegap0pt
\usepackage{amssymb}
\usepackage{enumerate}
\newtheorem{lemma}{Lemma}
\title{Geometry in the Furstenberg Conjecture}
\author{Yunping JIANG}
\date{Source: arXiv:2206.13569v2 (CC BY 4.0)}
\begin{document}
\maketitle

\noindent\emph{Source and licence.} Geometry in the Furstenberg Conjecture. Version arXiv:2206.13569v2 (CC BY 4.0), distributed by arXiv under the Creative Commons Attribution 4.0 International licence, http://creativecommons.org/licenses/by/4.0/, as recorded in the arXiv licence metadata for this exact version and shown on the abstract page https://arxiv.org/abs/2206.13569v2. Full licence notice: \texttt{licenses/arxiv-2206.13569-CC-BY-4.0.txt}.\par\smallskip

\noindent\textit{Editorial scope.} Counted unit: the article's lemma nt (source0.tex lines 662--672). The article's complete original proof of it is delivered with it (source0.tex lines 674--737, one \texttt{proof} environment of 64 lines). No mathematics is written, repaired or re-derived: the statement and the proof are the article's own lines, in the article's own notation, and no retained line is substituted or re-flowed.  No further source-local context is needed: every cross-reference of every retained line resolves inside the two delivered spans.  Nothing was omitted from any retained span, so the package carries no omission record.  No retained line contains a citation, so the wrapper carries no bibliography and no bibliography entry.  No retained line contains a figure environment, an image-inclusion command or drawing code, so no image is delivered and the package declares no asset dependency.  Editorial formatting only, in addition: an AMS \texttt{amsart} preamble that restates the article's own declarations and macros used by the retained lines, namely the environments \texttt{lemma} on separate counters, together with the standard packages those lines need.  The retained environments print in this document as Lemma 1 in order of appearance, whereas the article prints the same items with the numbering of its own full document.  The complete native source of the article as distributed in the arXiv e-print for this exact version is bundled unchanged as \texttt{sources/127499/source0.tex}.
\begin{lemma}~\label{nt} 
Suppose $p>1$ and $q>1$ are two relatively prime integers. Suppose $a\geq 1$ is an integer. 
Let 
$$
T_{n}=\#\Big(\{ aq^k \pmod{p^n}\;\;|\; k=0, 1, \cdots\}\Big)
$$
Then there is a constant $C_{1}=C_{1}(a,p,q)>0$ and a constant integer $n_{0}=n_{0}(a, p, q)$ such that 
$$
T_n=C_{1}p^n, \;\;\forall \; n\geq n_0.
$$
\end{lemma}

\begin{proof}
Let $a=bp^r$ with $\gcd(b,p)=1$. For every $n\geq r+1$, consider two integers $k>l$ 
such that $aq^k=aq^l \pmod{p^n}$. That is, 
$$
aq^l(q^{k-l}-1)= m p^n
$$
for some integer $m\geq 0$, 
which implies that 
$$
b q^l(q^{k-l}-1)=mp^{n-r}
$$
Since $b$ and $q$ are both relatively prime to $p$, we have $p^{n-r}|(q^{k-l}-1)$, that is,
$$
q^{k-l}=1\pmod{p^{n-r}}.
$$
Thus $T_{n}$ is the smallest positive integer such that $q^{T_{n}}=1\mod p^{n-r}$. 

Suppose $q^{T_n}=k_{n} p^{n-r}+1$ for some integer $k_{n}\geq 1$. 
For $n=r+1$, we have
$q^{T_{r+1}}=k_{r+1}p+1$.
Let us write $k_{r+1}=u_0p^s$, where $(u_0,p)=1$. Then
$$
q^{T_{r+1}}=u_0p^{s+1}+1=u_0p^{(r+s+1)-r}+1=q^{T_{r+s+1}},
$$
which implies
$$
q^{pT_{r+s+1}}=(u_0p^{s+1}+1)^p
$$
$$
=\sum_{j=0}^p\left(\begin{array}{l}p\cr                           j
\end{array}\right)(u_0p^{s+1})^{p-j} =(u_0p^{s+1})^p+p(u_0p^{s+1})^{p-1}+\ldots+p(u_0p^{s+1})+1
$$
$$
=u_1p^{s+2}+1.
$$
Note that
$u_1=u_0+p(\cdots)=u_0\pmod{p}$.
Since $(u_0,p)=1$, we have $(u_1,p)=1$ and thus
$$
u_1p^{s+2}+1=q^{T_{r+s+2}}.
$$
Therefore,
$$
pT_{r+s+1}=T_{r+s+2}.
$$
Inductively we get
$$
T_{r+s+1+m}=p^mT_{r+s+1}, \;\;\forall \; m\geq 0.
$$
Setting $n_0=r+s+1$ we have
$$
T_{n_0+m}=p^mT_{n_0}, \;\;\forall \; m\geq 0.
$$
That is,
$$
\frac{T_{n_0+m}}{p^{n_0+m}}=\frac{T_{n_0}}{p^{n_0}},\;\;\forall m\geq 0.
$$
Therefore, let
$$
C_{1}=\frac{T_{n_0}}{p^{n_0}}.
$$
Then
\[T_n=C_{1}p^n, \quad\forall n\geq n_0.\]
\end{proof}

\end{document}
